\BeleskeID{07-s030}
% element: 07-s030:d01
Sekvencijalno izvršavanje odnosi se na naredbe unutar jednog procesa. Različiti procesi u arhitekturi predstavljaju konkurentne aktivnosti. Lokalna promenljiva i signal imaju različitu semantiku dodele: promenljiva se menja odmah, dok signalna dodela zakazuje ažuriranje.

Lista osetljivosti navodi signale čiji događaji ponovo aktiviraju proces. Proces sa tom listom ponaša se kao da se na kraju izvršavanja suspenduje i čeka novi događaj na navedenim signalima. WAIT je drugačiji način određivanja mesta i uslova čekanja; proces sa eksplicitnom listom osetljivosti ne sme u svom telu imati WAIT naredbu. U sintaksnom obrascu prikazane su alternativne mogućnosti opisa.
